Every booster setting in this architecture was \textbf{copied from \texttt{OffsetModel} and never
tuned}. §22 required the capacity comparison to hold everything but the likelihood fixed, so the
probit objective was fitted at the \emph{softmax} objective's tuned settings. That kept the
comparison clean, and it means §11's negative result is a negative result about a borrowed
configuration rather than about the architecture at its best.

\textbf{One reading that looked like evidence, and was not.} §35's depth ablation is monotone from
$+0.002665$ at depth~1 to $-0.000814$ at depth~6 and stops there, which reads as ``the probit
objective wants deeper trees than production's depth~5''. It is not: the ablation reports probit
\emph{minus} softmax at each depth, and its own table has the softmax arm degrading with depth
(best at depth~1, $-0.004279$; $-0.003646$ at depth~5). Part of a narrowing gap was the other arm
getting worse. Measured directly on a fat fit, depth is flat and then harmful --- $-0.004200$,
$-0.004191$, $-0.004210$, $-0.003889$ for depths 5 to 8. \emph{A monotone difference between two
arms is not a monotone claim about either.}

\subsection*{The validation year is the whole problem}

The canonical profile scores 2015--2026 and trains from 2010, so only 2011--2014 are never scored.
That gives two options and no third: validate on 2014, which is leak-free but fits
\textbf{36{,}160} training rows where a canonical cell gets \textbf{121{,}000}; or validate on 2023
as §8 did, which has the rows and is a year the canonical run \emph{scores}, so the selection leaks
into that cell.

Both were run over one declared five-configuration grid. \texttt{utility\_scale} was fixed at the
declared 2.0 throughout --- a scale fitted on the training loss runs to its lower bound (§4), and
letting it move would make the sweep a comparison of scales. Scoring goes through the model's own
\texttt{predict\_proba}, so the loss cannot differ from the one the architecture reports.

\begin{center}\small
\begin{tabular}{lrrrrr}
\toprule
configuration & 2014 (36k) & 2022 (112k) & 2023 (121k) & fat mean & 2022/23 swing \\
\midrule
baseline (mcw 20, $\lambda$ 10) & $-0.000314$ & $-0.004371$ & $-0.004200$ & $-0.004285$ & $0.000171$ \\
\texttt{reg\_lambda=5.0} & $-0.000316$ & $-0.004514$ & $-0.003802$ & $-0.004158$ & $0.000712$ \\
\texttt{reg\_lambda=2.0} & \bad{$-0.000451$ \tiny(1st)} & \bad{$-0.005092$ \tiny(1st)} & \bad{$-0.003558$ \tiny(5th)} & $-0.004325$ & \bad{$0.001534$} \\
\texttt{gamma=0.0} & $-0.000288$ & $-0.004383$ & $-0.004186$ & $-0.004284$ & $0.000197$ \\
\rowcolor{rowshade}\texttt{min\_child\_weight=10} & \bad{$+0.000653$ \tiny(5th)} & $-0.004628$ & \good{$-0.004836$ \tiny(1st)} & \good{$-0.004732$} & $0.000208$ \\
\bottomrule
\end{tabular}
\end{center}

\textbf{The two fat years anti-rank the grid: Spearman $-0.4$, and 2014 against 2023 is $-0.9$.}
Every candidate is both first and last depending on which year is asked. The mechanism is
training-set size and it points the way it should --- a looser per-leaf minimum only pays when there
is data to fill the leaves --- which makes the leak-free year not merely weak but \emph{actively
misleading}: its whole signal is $-0.0003$ against the fat years' $-0.0043$, a factor of fourteen,
and its ranking is inverted. Selecting on 2014 would have chosen the configuration that comes last
on the fits that matter.

\subsection*{The selection rule, stated before the canonical run}

One year's argmax is the error §8 named and then half-committed. Two years allow a rule an argmax
does not:

\begin{quote}
Take the configuration that (a) improves on the baseline in \textbf{both} fat validation years and
(b) swings no more between them than the baseline itself does.
\end{quote}

Clause (b) disqualifies the largest average gain. \texttt{reg\_lambda=2.0} has the second-best fat
mean and buys it with a swing of $0.001534$ --- \textbf{nine times} the baseline's $0.000171$. A
setting whose value depends that strongly on which season validates it has been shown to be
sensitive, not to generalise. \texttt{min\_child\_weight=10} gains $+0.000447$ on the fat mean and
swings $0.000208$, statistically the control's own.

\textbf{Selected: \texttt{min\_child\_weight} $20 \to 10$, and nothing else.} Depth stays at 5,
\texttt{reg\_lambda} at 10.0, \texttt{gamma} at 0.2, the step at §8's 0.0025, rounds at 400 --- more
rounds was worse on both fat years (2023: 600 $\to -0.004011$, 800 $\to -0.003595$).

\textbf{2022 and 2023 selected it, so both are excluded from the canonical comparison it is judged
by}, leaving nine scored years of eleven. That is the price of using fat validation years, and
paying it explicitly is cheaper than the alternative: a tuned number whose validation year sits
inside its own test set.

\subsection*{The canonical result}

60 cells, five seeds, 400 rounds. \textbf{Control first}: the borrowed settings against the
production Offset reproduces §11's published figure exactly --- year-mean $-0.001352$, se $0.000968$,
$t=-1.40$, 6/11 years --- so what follows comes from the same machinery.

\begin{center}\small
\begin{tabular}{lrrrr}
\toprule
comparison, paired at year level, 9 leak-free years & year-mean & se & $t$ & years \\
\midrule
tuned vs borrowed settings & \bad{$+0.000062$} & $0.000099$ & \bad{$+0.63$} & 5/9 \\
tuned vs link-only probit & $+0.000054$ & $0.001046$ & $+0.05$ & 3/9 \\
tuned vs production Offset & $-0.000836$ & $0.001097$ & $-0.76$ & 4/9 \\
\rowcolor{rowshade}\textbf{probit link, tuned mean held fixed} & \good{$\mathbf{-0.000351}$} & $0.000108$ & \good{$\mathbf{-3.24}$} & \good{\textbf{8/9}} \\
\bottomrule
\end{tabular}
\end{center}

\textbf{The tuning did nothing.} A 15\% improvement in the validation year's edge ($-0.004836$
against $-0.004200$) became $+0.000062$ at $t=+0.63$ across nine independent years, alternating sign
year to year. The grid was declared in advance, the selection rule was stated before the run, that
rule deliberately rejected the largest average gain in favour of a stability criterion, and the
selection years were excluded from the judgement. Every precaution was taken and the gain still did
not survive. \emph{One validation year moved 15\% and nine test years moved nothing}; a single-year
delta at this effect size is not evidence about a hyperparameter, and this is the measurement that
says so rather than an argument that it might be.

\textbf{What the run did establish is the strongest number in the programme's record.} Scoring the
\emph{identical} boosted margins through the probit link and through the softmax link --- mean
function fixed, only the shock distribution changing --- gives $-0.000351$ at $\mathbf{t=-3.24}$ with
8 of 9 years favourable. §11 reported this contrast at \emph{seed} level, where five deltas over the
same races measure determinism rather than generalisation. At year level it is the first $t$ below
$-2$ in twenty candidates.

It is still not a promotion case, for §11's own reason: the contrast needs the boosted margins, which
cost two and a half hours of fitting and are worth nothing themselves --- \texttt{tuned vs link-only
probit} is $+0.000054$ at $t=+0.05$, so the free link over the incumbent's utilities is as good as
the expensive mean function beneath it. The link pays; the architecture that made it measurable is
not what should ship. \textbf{Recommendation unchanged: keep the softmax booster, and keep
\texttt{probit\_offset} as the published link.}
